%======================================================================
\section{Certified bounds}\label{sec:certified}
%======================================================================

\subsection{Method}\label{ssec:cert-method}

Every statement of this section is computer-assisted, and this paragraph states the conventions once for the whole paper. Each
proof names the function that is certified, the lemmas it uses and the arithmetic. Appendix~\ref{app:anc} gives, for each
statement, the script and its input, the command that re-runs it, the run time and the decisive output line; the SHA-256 hashes of
the inputs are listed there or in the file \texttt{SHA256SUMS} of the ancillary files. \emph{Exact} computations use integers and
rationals (FLINT \cite{FLINT}); \emph{ball} computations use Arb midpoint--radius arithmetic with outward rounding \cite{Joh17}, and a
sign is accepted only when the ball excludes~$0$. Every input is an exact rational parameter file, or a certified enclosure written by
another listed verifier, and its hash is recomputed. Printed upper bounds are rounded up and lower bounds down. Statements labelled ``Numerical observation'' are computed but
not certified, and are never used in a proof. The common lemmas and methods (positivity certificates, validated Bessel functions, the
exact members) are in Appendix~\ref{app:certification}. Theorem~\ref{thm:kappa}, Propositions~\ref{thm:kappaOPS}
and~\ref{prop:lowheight} and the upper bounds in Proposition~\ref{thm:notcritical} use only weak duality (Lemma~\ref{lem:weak-duality}) and
the exact membership of an explicit function in $\Cone$ or $\ConeOPS$: Proposition~\ref{prop:lowheight} bounds masses, and the others bound
slacks from above. Proposition~\ref{prop:lowerbound} and the lower bounds in Proposition~\ref{thm:notcritical} use explicit admissible pairs
(Lemma~\ref{lem:herglotz}). Propositions~\ref{thm:finiteJ} and~\ref{prop:exactcone} concern the exact family. No statement about the zeros of $\zeta$ is used.

\subsection{The slack}\label{ssec:kappa}

\begin{theorem}[A certified upper bound for the slack]\label{thm:kappa}
$\kstar\le\kOPS\le1.4291572\cdot10^{-1060}$. More precisely, for each $J$ in Table~\ref{tab:kappa-ladder} the exact member
$F_J=\Xi^2P_J(t^2)$ of Proposition~\ref{prop:E} lies in $\ConeOPS\subseteq\Cone$ and satisfies $\Arch(F_J)/\int F_J\le\kappa_J$, with
$\kappa_J$ as in the table; in particular $\kappa_{111}=1.4291572\cdot10^{-1060}$.
\end{theorem}

\begin{table}[ht]
\centering
\caption{Certified upper bounds $\kappa_J$ for $\kstar$ and $\kOPS$ from the exact members $F_J$ (Theorem~\ref{thm:kappa} and
Proposition~\ref{prop:exactcone}); $n_{J+1}$ is the first prime power not killed by $P_J$, and $x_0$ the start of the far field in the proof
of Proposition~\ref{prop:exactcone}.}\label{tab:kappa-ladder}
\begin{tabular}{rrrrr}
\toprule
$J$ & $n_{J+1}$ & $x_0$ & $\Arch(F_J)$ (rounded up) & $\kappa_J$ (rounded up)\\
\midrule
10 & 17 & 16.45 & $4.2125564\cdot10^{-19}$ & $1.5107258\cdot10^{-19}$\\
60 & 211 & 199.4 & $2.9669313\cdot10^{-417}$ & $1.0640143\cdot10^{-417}$\\
61 & 223 & 211.3 & $6.5071641\cdot10^{-442}$ & $2.3336286\cdot10^{-442}$\\
110 & 479 & 467.35 & $5.5796840\cdot10^{-1042}$ & $2.0010115\cdot10^{-1042}$\\
111 & 487 & 479.3 & $3.9851074\cdot10^{-1060}$ & $1.4291572\cdot10^{-1060}$\\
\bottomrule
\end{tabular}
\end{table}

\begin{proof}[Proof (computer-assisted; ball arithmetic; Appendix~\ref{app:anc}, Table~\ref{tab:anc-index}, first row)]
By Proposition~\ref{prop:exactcone} below, $F_J\in\ConeOPS\subseteq\Cone$, and $\Arch(F_J)$ and $\int F_J=\FT F_J(0)>0$ are enclosed in balls.
Since $F_J/\int F_J\in\ConeOPS$ has integral $1$, $\kstar\le\kOPS\le\Arch(F_J)/\int F_J$, and $\kappa_J$ is the upper end of the resulting
ball, rounded up.
\end{proof}

Every $F_J$ vanishes at every non-trivial zero of $\zeta$, on or off the line, so, by Proposition~\ref{prop:E}, $\Arch(F_J)$ is a sum over the
prime powers from $n_{J+1}$ on; its first term is of the size of $\Psi(\xi_{n_{J+1}})\approx C_\infty n_{J+1}^4e^{-2\pi n_{J+1}}$
(Proposition~\ref{thm:bessel}) times a cofactor. Table~\ref{tab:kappa-ladder} gives $-\log\kappa_J/n_{J+1}=2.55$, $4.55$, $4.56$, $5.01$, $5.01$
at $J=10$, $60$, $61$, $110$, $111$, against $2\pi\approx6.28$ for the factor $e^{-2\pi n_{J+1}}$ alone. So the number of certified digits is
governed by $n_{J+1}$, and it grows with $J$ for as long as positivity can be certified. As a cross-check, the cushioned
certificates $\Xi^2(H+\eps e^{-\pi t^2})$ of an earlier version of this paper, with $H$ a rational approximation of $P_J(t^2)$ and a
separate verifier (ancillary directory \nolinkurl{kappa/}), give $\kstar\le1.1508391\cdot10^{-906}$ at $J=100$, and at $J=60$ the bound
$1.0640143\cdot10^{-417}$, which agrees with $\kappa_{60}$ in Table~\ref{tab:kappa-ladder} to eight digits.

\begin{corollary}[Conductor form; infeasible deformations]\label{cor:qmin}
(a) No admissible pair has a zero measure $\mu\ge\lambda\dd t$ with $\lambda>1.4291572\cdot10^{-1060}$.
(b) If $\log q<-2\pi\cdot1.4291572\cdot10^{-1060}$, in particular if $q\le1-8.98\cdot10^{-1060}$, the data $\Arch_q$ admit
no admissible pair. Equivalently, replacing $\log\pi$ by $\log\pi+\eta$ in $\Oinf$ leaves no admissible pair once
$\eta>8.9796596\cdot10^{-1060}$. Hence $\qmin=e^{-2\pi\kstar}\ge1-8.98\cdot10^{-1060}$, and the classical conductor bound optimised over
$\ConeOPS$, $e^{-2\pi\kOPS}$, is at least $1-8.98\cdot10^{-1060}$ too. Under RH, $\qmin\le1$, so then $1-8.98\cdot10^{-1060}\le\qmin\le1$.
\end{corollary}

\begin{proof}
(a) is Lemma~\ref{lem:weak-duality} applied to $F_{111}/\int F_{111}$. (b) For that function,
$\Arch_q=\Arch+(\log q/2\pi)\int$ is negative when $\log q<-2\pi\kappa_{111}$, which excludes admissible pairs by
Lemma~\ref{lem:weak-duality}; replacing $\log\pi$ by $\log\pi+\eta$ is the case $\log q=-\eta$, and
$2\pi\kappa_{111}\le8.9796596\cdot10^{-1060}$. Then $\qmin\ge e^{-8.9796596\cdot10^{-1060}}\ge1-8.98\cdot10^{-1060}$
(Proposition~\ref{prop:conductor}), and $e^{-2\pi\kOPS}$ obeys the same bound because $F_{111}\in\ConeOPS$. Under RH, $\pz\in\Kad$, so
$\kstar\ge0$ (Theorems~\ref{thm:logic} and~\ref{thm:duality}).
\end{proof}

The same certificate also bounds how far the zero measure of any positive solution can stray from the zeros of $\zeta$.

\begin{corollary}[Near-criticality forces near-rigidity]\label{cor:nearrigid}
Every admissible pair $(\mu,\nu)$ satisfies $\int\Xi^2\dd\mu\le3.99\cdot10^{-1060}$. Consequently $\mu(I)\le3.99\cdot10^{-1060}/\min_I\Xi^2$ for
every compact interval $I\subset\RR\setminus\Zz$.
\end{corollary}

\begin{proof}[Proof (computer-assisted; ball arithmetic; Appendix~\ref{app:anc})]
By Proposition~\ref{prop:E}, $P_{111}(0)=1$, and by Proposition~\ref{thm:finiteJ} every coefficient of $P_{111}$ is positive. So $H_{111}\ge1$ on $\RR$ and
$F_{111}=\Xi^2H_{111}\ge\Xi^2$. As $F_{111}\in\Cone$ (Proposition~\ref{prop:exactcone}), Lemma~\ref{lem:weak-duality} gives
$\int\Xi^2\dd\mu\le\int F_{111}\dd\mu\le\Arch(F_{111})\le3.98510732132\cdot10^{-1060}$, the upper end of the enclosure of
$\Arch(F_{111})$ in Proposition~\ref{prop:exactcone}. On $I$, $\Xi^2\ge\min_I\Xi^2>0$.
\end{proof}

So every positive solution of the explicit formula, whether or not RH holds, has its zero measure concentrated on the zeros of $\zeta$, up to
$3.99\cdot10^{-1060}$ measured against $\Xi^2$. As $\Xi(t)^2$ decays like $e^{-\pi|t|/2}$, the bound is informative only at low height.
Every $F_J$ of Table~\ref{tab:kappa-ladder} satisfies $F_J\ge\Xi^2$ in the same way.

\subsection{The classical cone and low-height rigidity}

Let $f_k(u):=L_k^{(-1/2)}(2\pi u^2)e^{-\pi u^2}$; these are even Hermite functions, with $\widehat{f_k}=(-1)^kf_k$. They form the Laguerre basis
that Cohn and Gon\c{c}alves used for sign uncertainty \cite{CG19}.

\begin{proposition}[The classical cone]\label{thm:kappaOPS}
$\kstar\le\kOPS\le9.9462\cdot10^{-41}$. Explicitly, with $s=13$, $K=128$, explicit rationals $c_0,\dots,c_K$ and
$\eps=10^{-56}$, the function
\[
\Frep(t):=\sum_{k\le K}c_kf_k(t/s)+\eps\bigl(e^{-\pi t^2/(16s^2)}+e^{-16\pi t^2/s^2}\bigr)
\]
satisfies $\Frep>0$ and $\FT F_{\mathrm{rep}}>0$ on all of $\RR$, and $\Arch(\Frep)/\int\Frep\le9.9461827002\cdot10^{-41}$.
Consequently the classical explicit-formula conductor bound in degree~1, optimised over $\ConeOPS$, is at
least $1-6.25\cdot10^{-40}$.
\end{proposition}

\begin{proof}[Proof (computer-assisted; exact Sturm sequences and ball arithmetic; Appendix~\ref{app:anc})]
Write $X=2\pi t^2/s^2$, $Y=2\pi s^2\xi^2$ and $F:=\sum_{k\le K}c_kf_k(t/s)$. Then $F=e^{-X/2}P(X)$ and $\hF=s\,e^{-Y/2}Q(Y)$ with
$P,Q\in\QQ[X]$ of degree $128$. Exact Sturm sequences show that $P$ and $Q$ have no root in $(0,\infty)$, and $P(0)$,
$Q(0)$ and the leading coefficients are positive; hence $F>0$ and $\hF>0$ on $\RR$. The two Gaussians of the repair term are
positive with positive Fourier transforms, so $\Frep\ge F$ and $\FT F_{\mathrm{rep}}\ge\hF$, and $\Frep\in\ConeOPS\subseteq\Cone$. (Here
the repair term is not needed for positivity, since $P$ and $Q$ have no positive roots; its contribution to $\Arch(\Frep)$ is
included in the computed value.) Then $\Arch(\Frep)$ is one rigorous Arb integral of $\Frep$ against the digamma weight on a finite
interval, plus an explicit tail bound (Lemma~\ref{lem:digamma}), and $\int\Frep=\FT F_{\mathrm{rep}}(0)$ is enclosed exactly up to a ball. Finally
$2\pi\cdot9.9461827002\cdot10^{-41}\le6.2494\cdot10^{-40}$ and $e^{-a}\ge1-a$.
\end{proof}

Like $F_{111}$ (Theorem~\ref{thm:kappa}), this function lies in $\ConeOPS$ and needs no gap at $\xi_2$, so it is directly comparable with
Odlyzko's tables \cite[Table~4]{Odl90}; unlike $F_{111}$, it uses no arithmetic input. A certificate of the
same kind bounds the mass that \emph{any} admissible pair can put away from the low zeros of $\zeta$ and away from the prime powers.

\begin{proposition}[Low-height rigidity of every admissible pair]\label{prop:lowheight}
There is an explicit $F_{80}\in\ConeOPS$ of the form of Proposition~\ref{thm:kappaOPS}, with $K=80$, $s=12$ and $\eps=10^{-30}$,
such that $\Arch(F_{80})\le7.72525424595\cdot10^{-25}$ and $\int F_{80}\ge0.9999999999999999444$. Consequently every admissible pair
$(\mu,\nu)$ satisfies the mass bounds of Table~\ref{tab:windows}.
\end{proposition}

\begin{table}[ht]
\centering\small
\caption{Mass bounds valid for every admissible pair (Proposition~\ref{prop:lowheight}). The intervals $I_k$ are the gaps between the
first six positive zero ordinates of $\zeta$, shrunk by $0.1$ at both ends and rounded inward to ten decimals; the prime-side windows
are $[\xin a,\xin b]$.}\label{tab:windows}
\begin{tabular}{@{}lll@{\qquad}lll@{}}
\toprule
\multicolumn{3}{@{}l}{Zero side} & \multicolumn{3}{l@{}}{Prime side}\\
\cmidrule(r){1-3}\cmidrule(l){4-6}
$k$ & $I_k$ & bound for $\mu(I_k)$ & $[a,b]$ & & bound for $\nu$ \\
\midrule
1 & $[14.2347251418,20.9220396387]$ & $2.75\cdot10^{-17}$ & $[2.05,2.95]$ & & $6.81\cdot10^{-19}$\\
2 & $[21.1220396388,24.9108575801]$ & $1.91\cdot10^{-16}$ & $[5.05,6.95]$ & & $6.94\cdot10^{-14}$\\
3 & $[25.1108575802,30.3248761258]$ & $4.81\cdot10^{-15}$ & $[7.05,7.95]$ & & $2.72\cdot10^{-12}$\\
4 & $[30.5248761259,32.8350615877]$ & $2.28\cdot10^{-14}$ & $[10.05,10.95]$ & & $5.95\cdot10^{-10}$\\
5 & $[33.0350615878,37.4861781588]$ & $2.39\cdot10^{-13}$ & $[12.05,12.95]$ & & $3.80\cdot10^{-8}$\\
 & & & $[14,15]$ & & $1.82\cdot10^{-8}$\\
 & & & atom at $\xin{2.5}$ & & $1.53\cdot10^{-21}$\\
 & & & atom at $\xin6$ & & $6.37\cdot10^{-18}$\\
\bottomrule
\end{tabular}
\end{table}

No hypothesis on the zeros is used: whether or not RH holds, every admissible pair puts at most these masses into windows that avoid
the low zeros of $\zeta$ and the low prime powers. Positivity in the explicit formula has long been used to locate low zeros
\cite[p.~132]{Odl90}, \cite{Mil02,Bob15}; here the positivity of both measures takes the place of GRH.

\begin{proof}[Proof (computer-assisted; exact arithmetic and ball arithmetic; Appendix~\ref{app:anc})]
As in the proof of Proposition~\ref{thm:kappaOPS}, exact Sturm sequences give $F_{80}\in\ConeOPS$, and ball arithmetic gives the bounds for
$\Arch(F_{80})$ and $\int F_{80}$. For
$F\in\Cone$, $\mu$ even and $I\subset(0,\infty)$, Lemma~\ref{lem:weak-duality} gives $2\min_IF\cdot\mu(I)\le\Arch(F)$, and
$\min_I\hF\cdot\nu(I)\le\pi\Arch(F)$. Lower bounds $F_{80}\ge m$ on $40$ intervals in $t$ and $\FT F_{80}\ge m$ on $19$ windows in $\xi$ are
proved in exact arithmetic: $e^{X/2}$ is bounded by a rational Taylor polynomial with Lagrange remainder, and positivity of the resulting
rational polynomials is shown by Descartes' rule after a M\"obius map, with bisection (Lemma~\ref{lem:posc}). The values $\FT F_{80}(\xin n)$
at the two atoms are enclosed in ball arithmetic. With the certified upper bound for $\Arch(F_{80})$, the two inequalities give the
bounds of Table~\ref{tab:windows}, rounded up.
\end{proof}

\subsection{The exact family at finite \texorpdfstring{$J$}{J}}

\begin{proposition}[Certified finite-$J$ instances]\label{thm:finiteJ}
For $J\in\{10,60,61,110,111\}$ the Hermite matrix $M_J$ is non-singular, every coefficient of the exact $P_J$ is positive, and
$\max_{1\le k\le2J}((2k)!\,p_k^{(J)})^{1/2k}\le0.71636$, that is, $0<p_k^{(J)}\le0.71636^{2k}/(2k)!$. Hence $H_J\ge1$ on $\RR$ and
$\lvert H_J(t)\rvert\le\cosh(0.71636\lvert t\rvert)$ on $\CC$.
\end{proposition}

\begin{proof}[Proof (computer-assisted; ball arithmetic; Appendix~\ref{app:anc})]
Through Proposition~\ref{thm:bessel}, the moments $(-1)^k(2\pi)^{-2k}\Psi^{(2k+\varepsilon)}(\xi_n)$, with $\varepsilon\in\{0,1\}$, are
enclosed rigorously; rows and columns are scaled by exact powers of two, and the ball system is solved with Arb's preconditioned
solver. A successful solve proves that every matrix in the ball, in particular $M_J$, is non-singular, and the output balls contain the
exact $p^{(J)}_k$; positivity holds because every ball is positive. The certified bounds for $\max_k((2k)!p_k)^{1/2k}$ are $0.697078$
($J=10$), $0.716207$, $0.716214$, $0.716351$ and $0.716352$; the bounds on $H_J$ follow from Proposition~\ref{prop:coeff} with $C=1$.
Each enclosure is compared with a second enclosure computed by a different implementation, and every pair of balls overlaps.
\end{proof}

Uniform versions of these facts are the hypotheses \hyp{N}, \hyp{G$_a$} and \hyp{Z$_\infty$} of Theorem~\ref{thm:criterion}, through
Proposition~\ref{prop:coeff}. At the same values of $J$ the exact members lie in the cone, without approximation or cushion; the method is in
Appendix~\ref{app:exact}.

\begin{proposition}[Exact members in the cone]\label{prop:exactcone}
For $J\in\{10,60,61,110,111\}$ the exact member $F_J=\Xi^2P_J(t^2)$ of Proposition~\ref{prop:E} satisfies:
\textup{(a)} $F_J\ge\Xi^2$ on $\RR$;
\textup{(b)} $\FT F_J(\xi)>0$ for every $\xi\ge0$ except at $\xi_{n_1},\dots,\xi_{n_J}$, where $\FT F_J$ has zeros of order exactly two.
Hence $F_J\in\ConeOPS\subseteq\Cone$. Moreover, \textup{(c)} $\Arch(F_J)=\frac1\pi\sum_{n\in\PP,\,n>n_J}\Lambda(n)n^{-1/2}\FT F_J(\xi_n)$ is at
most the value in Table~\ref{tab:kappa-ladder}, and $2.7884318\le\int F_J=\FT F_J(0)\le2.7884322$; at $J=111$,
$\Arch(F_{111})\in[3.98510732131,3.98510732132]\cdot10^{-1060}$.
\end{proposition}

\begin{proof}[Proof (computer-assisted; ball arithmetic; Appendices~\ref{app:exact} and~\ref{app:anc})]
(a) Since $P_J(0)=1$ (Proposition~\ref{prop:E}) and every coefficient of $P_J$ is positive (Proposition~\ref{thm:finiteJ}), $H_J\ge1$; and $F_J\in\Tc$ by
Section~\ref{ssec:bessel}. (b) The input is the ball
enclosure of $p^{(J)}$ of Proposition~\ref{thm:finiteJ}, which contains the exact coefficients. With $\FT F_J$ and its derivatives evaluated
through Proposition~\ref{thm:bessel} and Lemma~\ref{lem:abel}, $[0,\infty)$ minus the nodes is covered exactly by punctured neighbourhoods of the nodes, where
$\FT F_J(\xi_n)=\FT F_J'(\xi_n)=0$ by Proposition~\ref{prop:E} and a positive lower bound for $\FT F_J(\xi_n+\delta)/\delta^2$ is certified,
by Taylor pieces between them, and by the far field $x\ge x_0$, where Lemma~\ref{lem:abel} gives positivity. As $\FT F_J$ is even,
$\FT F_J\ge0$ on $\RR$. (c) By Corollary~\ref{cor:zerokilling} and Proposition~\ref{prop:E}, with no hypothesis on the zeros of $\zeta$,
$\Arch(F_J)$ is the stated sum. Its terms are non-negative by (b); those with $n\le400$ ($J\le61$) or $n\le800$ ($J\ge110$) are enclosed,
and the remaining terms sum to at most $6.2\cdot10^{-998}$, $2.9\cdot10^{-835}$, $9.3\cdot10^{-834}$, $1.3\cdot10^{-1760}$ and
$2.6\cdot10^{-1759}$ for $J=10$, $60$, $61$, $110$, $111$. Each certificate was reproduced through a second evaluation path, with
independently computed coefficient enclosures, an independent operator construction and the other validated $K_0,K_1$, and with the same
certificate logic.
\end{proof}

\begin{observation}[The exact family for $J\le110$]\label{obs:family}
The following were computed in high-precision arithmetic (at most $4600$ bits); the floating-point table is
\nolinkurl{computed/family_J1_110.tsv} in the ancillary files. For every $1\le J\le110$ the Hermite system is non-singular and
$p^{(J)}_{2J}>0$. All coefficients $p_k^{(J)}$ are positive for $J\in\{1,2\}$ and for $10\le J\le110$; for $3\le J\le9$ the coefficient
$p^{(J)}_{2J-1}$ is non-positive, and for $4\le J\le8$ so is $p^{(J)}_{2J-3}$. The quantity $a_J:=\max_{1\le k\le2J}((2k)!\,|p_k^{(J)}|)^{1/2k}$ (for $J\le110$ the maximum is attained at a positive
coefficient, so the table's values are the same if only positive coefficients are used) stays below $0.71$ for $J\le9$ and increases with $J$ for $10\le J\le110$, up to $0.7164$ at $J=110$ (attained near $k\approx20$--$24$); so \eqref{eq:POSa} holds
numerically with $a=0.7164$ and $C=1$ for $10\le J\le110$. Dense scans at $J=20,30,40,50,60$ ($x\in[2,1.15n_J+5]$, step $0.05$, and up to
$x=425$ at $J=60$) found no sign change of $\FT F_J$ off the nodes; for $J\in\{10,60,61,110,111\}$ this is certified
(Proposition~\ref{prop:exactcone}).
\end{observation}

\subsection{A lower bound for the slack, and other archimedean data}\label{ssec:otherdata}

Lower bounds for slacks come from explicit admissible pairs, which are feasible points of the dual linear programme; Cohn and
Triantafillou used dual feasible points in the same way to prove that linear programming bounds for sphere packing are not sharp
\cite{CT22}. We build them for the data of $\zeta$ with a conductor $q$, that is, for
$\Arch_q=\Arch+\frac{\log q}{2\pi}\int$ (Section~\ref{ssec:criticality}), and, as a control, for the archimedean data of $\QQ(\sqrt5)$
and $\QQ(\sqrt{-3})$. For $\mathfrak g\in\{\GR,\GR^2,\GC\}$ let $(k_j)_{j\le d}$ be $(0)$, $(0,0)$ and $(0,1)$ respectively, put
\begin{gather*}
\Omega_{\mathfrak g}(t):=\frac1{2\pi}\sum_{j\le d}\Bigl(\Re\psi\bigl(\tfrac{\frac12+k_j+it}{2}\bigr)-\log\pi\Bigr),\\
\Arch_{\mathfrak g,q}(F):=F\bigl(\tfrac i2\bigr)+F\bigl(-\tfrac i2\bigr)+\int_\RR F(t)\Bigl(\Omega_{\mathfrak g}(t)+\frac{\log q}{2\pi}\Bigr)\dd t,
\end{gather*}
so that $\Omega_{\GR}=\Oinf$ and $\Arch_{\GR,q}=\Arch_q$. Admissible pairs, $\Cone$ and $\kstar$ are defined for $\Arch_{\mathfrak g,q}$ exactly as for
$\zeta$, with the same test class and the same gap $\xi_2$ (Definition~\ref{def:admissible}); Lemma~\ref{lem:weak-duality} holds with the
same proof, and $\kstar(\Arch_{\mathfrak g,q})=\kstar(\Arch_{\mathfrak g,1})+\frac{\log q}{2\pi}$. More generally, for a gap $g>0$ we write $\Cone_g$ and $\kstar_g$ for the cone and the slack defined with $[g,\infty)$ in place of
$[\xi_2,\infty)$, both for the prime measures and for the condition $\hF\ge0$; Lemma~\ref{lem:weak-duality} holds with the same proof, and
$\ConeOPS\subseteq\Cone_g$ for every $g$.

\begin{lemma}[Herglotz form of an admissible pair]\label{lem:herglotz}
Let $\tilde\nu\ge0$ be a measure on $[2,\infty)$ such that $\sigma:=\tilde\nu-\one_{[2,\infty)}(x)\dd x$ satisfies
$\int x^{-1/2}\dd\lvert\sigma\rvert(x)<\infty$, and put $\mathcal L(s):=\int x^{-s}\dd\tilde\nu(x)$ for $\Re s>1$. Then $\mathcal L$ extends
continuously to $\Re s\ge\frac12$, $s\ne1$, as $2^{1-s}/(s-1)+\int x^{-s}\dd\sigma$. Let
$\mu(t):=\Omega_{\mathfrak g}(t)+\frac{\log q}{2\pi}-\frac1\pi\Re\mathcal L(\frac12+it)$, and let $\nu$ be the image of $x^{-1/2}\tilde\nu$ under
$x\mapsto\log x/2\pi$. Then $\Arch_{\mathfrak g,q}(F)=\int F\mu\dd t+\frac1\pi\int\FT F\dd\nu$ for every $F\in\Tc$. Hence $(\mu\dd t,\nu)$ is an
admissible pair for $\Arch_{\mathfrak g,q}$ if and only if $\mu\ge0$.
\end{lemma}

\begin{proof}
Let $F\in\Tc_\delta$ and $\frac12<y<\frac12+\delta$. Since $\FT F(\log x/2\pi)=\int F(t)x^{-it}\dd t$, shifting the contour
to $\Im t=-y$ and using Fubini (note $\int x^{-1/2-y}\dd\tilde\nu<\infty$) gives
$\frac1\pi\int\FT F\dd\nu=\frac1\pi\int_{\Im w=-y}F(w)\mathcal L(\frac12+iw)\dd w$. The integrand is meromorphic in
$-y\le\Im w\le0$ with one simple pole, at $w=-i/2$, where $\mathcal L(\frac12+iw)\sim1/(i(w+\frac i2))$; $F$ decays like
$(1+\lvert w\rvert)^{-2}$ and $\mathcal L$ is bounded away from the pole. Moving the contour to $\Im w=-\eta$, $0<\eta<\frac12$, picks up
$2\pi i\cdot F(-\frac i2)/i$. As $\eta\downarrow0$ the integrals over $\Im w=-\eta$ converge to the integral over $\RR$, by dominated
convergence: $\mathcal L(\frac12+iw)$ is bounded on $-\frac14\le\Im w\le0$ away from the pole and continuous up to $\Im w=0$. So
$\frac1\pi\int\FT F\dd\nu=\frac1\pi\int_\RR F(t)\mathcal L(\frac12+it)\dd t+2F(-\frac i2)$. As $F$ is even and real and
$\mathcal L(\frac12-it)=\overline{\mathcal L(\frac12+it)}$, the integral equals $\frac1\pi\int F\,\Re\mathcal L(\frac12+it)\dd t$, and
$2F(-\frac i2)=F(\frac i2)+F(-\frac i2)$. Adding $\int F\mu$ gives $\Arch_{\mathfrak g,q}(F)$. The measure $\nu$ is positive
and carried by $[\xi_2,\infty)$, and $\mu$ is even; so admissibility is equivalent to $\mu\ge0$.
\end{proof}

Lemma~\ref{lem:herglotz} is a representation of Herglotz--Nevanlinna type: the zero density is read off from the boundary values of
$\Re\mathcal L$ on the critical line, much as Vedana's classification of summation formulas on a strip describes the continuous part of
the Fourier-side measure through a generalised Nevanlinna representation of a generating function \cite{Ved26}.

We use measures $\tilde\nu$ equal to $\one_{[2,\infty)}\dd x$ minus densities on finitely many cells, plus finitely many atoms and polynomial
bumps, minus tails $a_mx^{-m}\dd x$ ($m\ge1$) beyond some $X$; then $\mathcal L$ and $\mu$ are explicit. The verifier checks $\tilde\nu\ge0$ and
$\supp\tilde\nu\subseteq[2,\infty)$, and proves $\mu>0$ on a range $[0,T]$ by mean-value steps $\mu(t)\ge\mu(t_0)-h\sup\lvert\mu'\rvert$, with
$\mu'$ enclosed in ball arithmetic on each whole step. On $[T,\infty)$ it uses monotone majorants: $\Omega_{\mathfrak g}$ is increasing in $t\ge0$,
while the contributions of the cells, bumps and tails to $\mu$ are bounded by explicit constants times the decreasing functions
$(\frac14+t^2)^{-1/2}$, $\min(1,(k/\beta t)^k)$ (for bumps of half-width $\beta$ built from $k$ uniform variables) and
$((m-\frac12)^2+t^2)^{-1/2}$. The tail coefficients are checked to satisfy $a_m\ge0$, so that $1-\sum_ma_mx^{-m}$ is increasing and
$\tilde\nu\ge0$ needs to be checked only at $x=X$.

\begin{proposition}[An unconditional lower bound for the slack]\label{prop:lowerbound}
For $\zeta$'s data with conductor $q=e^{0.02}$ there is an admissible pair of the form of Lemma~\ref{lem:herglotz} with
$\mu\ge4.71949\cdot10^{-4}$ on $\RR$. Hence $\qmin\le e^{0.02}$ and
\[
\kstar\ge4.71949\cdot10^{-4}-\frac{0.02}{2\pi}\ge-2.7112\cdot10^{-3}.
\]
\end{proposition}

\begin{proof}[Proof (computer-assisted; ball arithmetic; Appendix~\ref{app:anc})]
The pair is certified as described above, with mean-value steps on $[0,2000]$ and monotone majorants beyond; Lemma~\ref{lem:herglotz} gives
admissibility for $\Arch_q$, so $\qmin\le e^{0.02}$. For $F\in\Cone$ with $\int F=1$, Lemma~\ref{lem:weak-duality} gives
$\Arch_q(F)\ge\int F\mu\dd t\ge4.71949\cdot10^{-4}$; so $\kstar(\Arch_q)\ge4.71949\cdot10^{-4}$, and $\kstar=\kstar(\Arch_q)-0.02/(2\pi)$ by
Proposition~\ref{prop:conductor}. Finally $0.02/(2\pi)\le3.1831\cdot10^{-3}$.
\end{proof}

\begin{proposition}[The data of $\QQ(\sqrt5)$ and $\QQ(\sqrt{-3})$, also at their natural gaps]\label{thm:notcritical}
\textup{(a)} For the archimedean data $\Arch_{\GR^2,5}$ of $\QQ(\sqrt5)$ and every gap $g\in(0,\xin{4.04915}]$, in particular for $\zeta$'s gap
$\xi_2$ and for the natural gap $\xi_4$ of the field: $1.64\cdot10^{-5}\le\kstar_g\le4.83\cdot10^{-4}$. In particular some admissible pair has zero
measure $\ge1.64\cdot10^{-5}\dd t$, and the set of admissible pairs is not a singleton.
\textup{(b)} For the data $\Arch_{\GC,3}$ of $\QQ(\sqrt{-3})$ and every gap $g\in(0,\xin{3.011664}]$, in particular for $\xi_2$ and for the natural
gap $\xi_3$: $1.50\cdot10^{-4}\le\kstar_g\le9.99\cdot10^{-4}$, with the same consequences.
\end{proposition}

\begin{proof}[Proof (computer-assisted; exact Sturm sequences and ball arithmetic; Appendix~\ref{app:anc})]
(a) For the upper bound, an explicit $F$ of the form of Proposition~\ref{thm:kappaOPS}, a combination of $f_0(t/7),\dots,f_{80}(t/7)$ plus the two
Gaussians, has $\Arch_{\GR^2,1}(F)/\int F\le k_0:=-0.25566705789831547610$; positivity of $F$ and $\hF$ on $\RR$ is shown by exact Sturm sequences,
so $F\in\ConeOPS\subseteq\Cone_g$ for every $g$, and the functional is one Arb integral. Hence $\kstar_g(\Arch_{\GR^2,5})\le k_0+\log5/2\pi\le4.83\cdot10^{-4}$.
For the lower bound, a pair of the form of Lemma~\ref{lem:herglotz} for $\Arch_{\GR^2,q_0}$, $\log q_0=1.6093347792651136$, has $\tilde\nu\ge0$, and
$\tilde\nu$ vanishes on $[2,x_0)$ for some $x_0\ge4.049150$, a bound read off exactly from its parameters; its zero density satisfies
$\mu\ge6.59\cdot10^{-9}$ on $[0,800]$ and $\mu\ge1.79$ on $[800,\infty)$. So the pair is admissible for every gap $g\le\xin{x_0}$. Read at $q=5$, the
same prime measure gives the zero density $\mu+(\log5-\log q_0)/2\pi\ge1.64\cdot10^{-5}$, and Lemma~\ref{lem:weak-duality} for the gap $g$ gives
$\kstar_g\ge1.64\cdot10^{-5}$. Adding to $\tilde\nu$ a small atom $w\delta_{x_1}$, $x_1\ge x_0$, changes $\mu$ by $-(w/\pi)x_1^{-1/2}\cos(t\log x_1)$; for
small $w>0$ the result is a second admissible pair. The two certificates are combined in exact arithmetic, which gives
$1.6420748\cdot10^{-5}\le\kstar_g\le4.8294147\cdot10^{-4}$.
(b) Identical, with a dual certificate for $\GC$ ($K=64$, functions $f_k(t/9)$), also in $\ConeOPS$, and a pair of the form of
Lemma~\ref{lem:herglotz} at $\log q_0=1.0976698108720329$ whose prime measure vanishes below $x_0\ge3.011664$ and whose zero density satisfies the
certified floor $\mu\ge6.67\cdot10^{-8}$ on $\RR$. The lower bound is $(\log3-\log q_0)/2\pi+\inf\mu$; the floor is needed, since
$(\log3-\log q_0)/2\pi$ alone is slightly below $1.5\cdot10^{-4}$. The combination gives $1.5006672\cdot10^{-4}\le\kstar_g\le9.9854968\cdot10^{-4}$.
\end{proof}

\begin{remark}[The gap and the margins]\label{rem:gap}
The prime-ideal norms of $\QQ(\sqrt5)$ start at $4$, and those of $\QQ(\sqrt{-3})$ at $3$, so $\xi_4$ and $\xi_3$ are the gaps that match the
fields' own explicit formulas. Proposition~\ref{thm:notcritical} holds at these natural gaps as well as at $\xi_2$, so for these two fields the
comparison with $\zeta$ is like-for-like. Under GRH for the Dedekind zeta function, the field's own pair is admissible at its natural gap.
Since $\kstar_g>0$ there, these data are not critical, and their admissible pairs are not unique: the field's own pair, whose zero measure is
atomic, and the certified pair, whose zero measure is at least $c\dd t$ with $c>0$, are two distinct pairs. This indicates that
near-criticality is a property of $\zeta$'s data rather than of the method. These are two examples; we make no claim about $L$-functions in
general. The margins are small. At every gap $g\le\xin{4.04915}$ the certificates in (a) place the critical conductor of the data $\GR^2$ (the
infimum of the $q$ for which $\Arch_{\GR^2,q}$ has an admissible pair at the gap $g$) in $[4.9848,4.9995]$: the dual function excludes every
$q<4.9848$ by weak duality, and the pair is admissible at $q_0<4.9995$. So the conductor $5$ of $\QQ(\sqrt5)$ exceeds it by a factor between
$1+1.0\cdot10^{-4}$ and $1+3.1\cdot10^{-3}$. For $\GC$ and every gap $g\le\xin{3.011664}$ the same argument gives $[2.9812,2.9972]$, and the
conductor $3$ of $\QQ(\sqrt{-3})$ exceeds it by a factor between $1+9.4\cdot10^{-4}$ and $1+6.3\cdot10^{-3}$. The slack as a
function of the gap is treated in \cite{N3}.
\end{remark}
